\chapter{Review of applied linear algebra}
\label{ch:Lec0-linalg}

The following is a brief review of some basic concepts in linear algebra that will be used throughout the course. Many definitions and results are stated 
without proof. The reader is encouraged to consult standard references 
such as \cite{olver2006applied,boyd2018introduction,demmel1997applied,matrixcookbook}.


\section{Vectors, inner products, and bases}

Write $\R^d$ to denote the $d$-dimensional Euclidean space. 
We will use boldface letter to denote vectors in $\R^d$,
for example $\xb \in \R^d$ to denote the vector 
$\xb = (x_1, x_2, \ldots, x_d)$ with $x_j \in \R$ for $j=1,2,\ldots,d$
denoting the individual entries. The usual Euclidean i`nner product of 
two vectors $\xb, \yb \in \R^d$ is defined as
\begin{equation}\label{def:euclidean-inner-product}
\langle \xb, \yb \rangle := \sum_{j=1}^d x_j y_j.    
\end{equation}

\begin{definition}[Inner product]
    Any function $\iota: \R^d \times \R^d \to \R$ that satisfies 
    the following four axioms is called an inner product on $\R^d$:
    \begin{itemize}
        \item Symmetry: $\iota(\xb, \yb) = \iota(\yb, \xb)$ for all $\xb, \yb \in \R^d$
        \item Linearity: $\iota(\alpha \xb + \beta \yb, \zb) = \alpha \iota(\xb, \zb) + \beta \iota(\yb, \zb)$ for all $\xb, \yb, \zb \in \R^d$ and $\alpha, \beta \in \R$
        \item Positivity: $\iota(\xb, \xb) \ge 0$ for all $\xb \in \R^d$
        \item Definiteness: $\iota(\xb, \xb) = 0$ if and only if $\xb = 0$
    \end{itemize}
\end{definition}
The Euclidean inner product defined in  \eqref{def:euclidean-inner-product}
is indeed an inner product on $\R^d$ according to the above definition.
The proof is straightforward and is left as an exercise for the reader.
We use the bracket notation to distinguish this inner product from 
other inner products that we may encounter later.


\begin{definition}[Norm]
    We say that a function $\|\cdot\|: \R^d \to \R$ is a norm on $\R^d$ if it satisfies the following three axioms:
    \begin{itemize}
        \item Positivity: $\|\xb\| \ge 0$ for all $\xb \in \R^d$
        \item Definiteness: $\|\xb\| = 0$ if and only if $\xb = 0$
        \item Triangle inequality: $\|\xb + \yb\| \le \|\xb\| + \|\yb\|$ for all $\xb, \yb \in \R^d$
    \end{itemize}
\end{definition}
The Euclidean inner product induces the Euclidean norm on $\R^d$, 
also routinely called the $2$-norm or the $\ell^2$-norm, defined as
\begin{equation*}
    \|\xb\|_2 = \sqrt{\langle \xb, \xb \rangle} = \sqrt{\sum_{j=1}^d x_j^2}.
\end{equation*}

With the above notion of inner product and norm, we can define the angle between two vectors and the notion of orthogonality.
The angle between two vectors $\xb$ and $\yb$ is defined as
\begin{equation*}
   \theta(\xb, \yb) \equiv \theta := \arccos \left( \frac{\langle \xb, \yb \rangle}{\|\xb\|_2 \|\yb\|_2} \right),
\end{equation*}
Hence, we say $\xb$ and $\yb$ are orthogonal (with respect to 
the Euclidean inner product) if $\langle \xb, \yb \rangle = 0$,
since it implies that $\theta = \pi/2$.

We say that a set of vectors $\{\xb_1, \xb_2, \ldots, \xb_N\} \equiv 
\{ \xb_j \}_{j=1}^N \subset \R^d$ 
are {\it linearly independent} if the only solution to the equation
\begin{equation*}
    \sum_{j=1}^N \alpha_j \xb_j = 0,
\end{equation*}
is $\alpha_j = 0$ for all $j=1,2,\ldots,N$ otherwise we 
say that they are {\it linearly dependent}.
For a set of vectors $\{\xb_j\}_{j=1}^N$, we also define their {\it span} as
\begin{equation*}
    \Span\{\xb_j\}_{j=1}^N := \left\{ \sum_{j=1}^N \alpha_j \xb_j : \alpha_j \in \R, j=1,2,\ldots,N \right\} \subseteq \R^d.
\end{equation*}
For a subspace $V \subseteq \R^d$, we say that a set of vectors
$\{\xb_j\}_{j=1}^N$ span $V$ if $\Span\{\xb_j\}_{j=1}^N = V$.
We say that a set of vectors $\{\xb_j\}_{j=1}^N$ are a {\it basis} 
for $V$ if they are linearly independent and they span $V$. Bases are 
not necessarily unique, i.e., we can find two sets of linearly 
independent vectors that span the same subspace $V$. However,
the number of vectors in any basis for a subspace $V$ is unique and is called the {\it dimension} of $V$.
\begin{lemma}
    Every subspace $V \subseteq \R^d$ has a basis and any two bases for $V$ have the same number of elements. The number of elements in a basis for $V$ is called the dimension of $V$ and is denoted by $\dim(V)$.
\end{lemma}
Bases are also useful for representing elements of a subspace.

\begin{lemma}
If $\{ \xb_j\}_{j=1}^N$ form a basis for a subspace $V$ then every element 
$\vb \in V$ can be uniquely expressed as a linear combination of 
the $\xb_j$'s, i.e., there exist unique coefficients $\alpha_j \in \R$ such that
$\vb = \sum_{j=1}^N \alpha_j \xb_j.$    
\end{lemma}




We say  $\{\xb_j\}_{j=1}^N$ 
are {\it orthogonal} if $\langle \xb_i, \xb_j \rangle = 0$ for all $i \neq j$.
We further say they are {\it orthonormal} if they are orthogonal and $\|\xb_j\|_2 = 1$ for all $j=1,2,\ldots,N$. 

\begin{lemma}
    Every orthogonal set of non-zero vectors is linearly independent.
\end{lemma}

Given a linearly independent set of vectors $\{ \xb_j \}_{j=1}^N$, 
we can always construct an orthonormal set of vectors $\{\qb_j\}_{j=1}^N$
such that 
\begin{equation*}
    \Span\{\qb_j\}_{j=1}^N = \Span\{\xb_j\}_{j=1}^N.
\end{equation*}
The classic procedure for doing so is called the Gram-Schmidt process: 
\begin{algorithm}
\caption{Gram-Schmidt Process}
\begin{algorithmic}
    \State {\bf Input}: linearly independent set of vectors 
    $\{\xb_j\}_{j=1}^N \subseteq \R^d$.
    \State {\bf Output}: orthonormal set of vectors $\{\qb_j\}_{j=1}^N \subset \R^d$
    \vspace{2ex}
    \State Set $\qb_1 = \xb_1 / \|\xb_1\|_2$.
    \For{$j=2,3,\ldots,N$}
        \State 
        \begin{equation*}
            \qb_j = \frac{\xb_j - \sum_{i=1}^{j-1} \langle \xb_j, \qb_i \rangle \qb_i}{\left\| \xb_j - \sum_{i=1}^{j-1} \langle \xb_j, \qb_i \rangle \qb_i \right\|_2}.
        \end{equation*}
    \EndFor
\end{algorithmic}
\end{algorithm}


\section{Matrices}

For positive integers $m, n$ consider a function $f: \R^n \to \R^m$. 

\begin{definition}[Linear and affine functions]
We say $f: \R^n \to \R^m$ is linear if it satisfies the following two properties:
\begin{itemize}
    \item Additivity: $f(\xb + \yb) = f(\xb) + f(\yb)$ for all $\xb, \yb \in \R^n$
    \item Homogeneity: $f(\alpha \xb) = \alpha f(\xb)$ for all $\xb \in \R^n$ and $\alpha \in \R$
\end{itemize}
We say $g: \R^n \to \R^m$ is affine if there exists a linear function $f: \R^n \to \R^m$ and a vector $\vb \in \R^m$ such that $g(\xb) = f(\xb) + \vb$ for all $\xb \in \R^n$.
\end{definition}

Linear functions from $\R^n$ to $\R$ can be represented as matrix-vector products. More precisely, 
\begin{lemma}
Every linear function $f: \R^n \to \R$ is uniquely associated with 
   with a vector $\fb \in \R^n$ such that $f(\xb) = \langle \fb, \xb \rangle$ for all $\xb \in \R^n$. 
\end{lemma}
In other words, linear functions from $\R^n$ to $\R$ can be identified with (are isomorphic to) vectors in $\R^n$. Additionally, for $f: \R^n \to \R^m$ write $f(\xb) = (f_1(\xb), f_2(\xb), \ldots, f_m(\xb))$ where $f_j: \R^n \to \R$ for $j=1,2,\ldots,m$. Then $f$ is linear if and only if each $f_j$ is linear. Hence, we can represent a linear function $f: \R^n \to \R^m$ 
\begin{equation*}
    f(\xb) = \begin{bmatrix}
        f_1(\xb) \\
        f_2(\xb) \\
        \vdots \\
        f_m(\xb)
    \end{bmatrix} = 
    \begin{bmatrix}
        \langle \fb_1, \xb \rangle \\
        \langle \fb_2, \xb \rangle \\
        \vdots \\
        \langle \fb_m, \xb \rangle
    \end{bmatrix} 
\end{equation*}
Collecting the vectors $\fb_j$ into a list, as the rows of a table, we obtain 
the classic definition of a matrix $F$. 
\begin{definition}
    An $m \times n$ matrix $F$ is a rectangular array of real numbers with $m$ rows and $n$ columns. The entries in the $i$-th row and $j$-th column of $F$ are denoted by $f_{ij}$. The set of all matrices with $m$ rows and $n$ columns is denoted by $\R^{m \times n}$.
\end{definition}

Simply put, given a matrix $F \in \R^{m \times n}$, we think of  
\begin{equation*}
    F = \begin{bmatrix}
        f_{11} & f_{12} & \cdots & f_{1n} \\
        f_{21} & f_{22} & \cdots & f_{2n} \\
        \vdots & \vdots & \ddots & \vdots \\
        f_{m1} & f_{m2} & \cdots & f_{mn}
    \end{bmatrix} 
    \equiv 
    \begin{bmatrix}
        -  & \fb_{1:} &  - \\
        -  & \fb_{2:} &  - \\
         & \vdots &  \\
        -  & \fb_{m:} &  -
    \end{bmatrix} 
    \equiv 
    \begin{bmatrix}
        |        & |        &  |  &      | \\
        \fb_{:1} & \fb_{:2} & \cdots & \fb_{:n} \\
                |        & |        &  |  &      |   
    \end{bmatrix}
\end{equation*}
where we also introduced the \texttt{Python}-style notation $\fb_{j:}$ to 
denote the $j$-th row of $F$ and $\fb_{:j}$ to denote the $j$-th column of $F$.
We will also use the standard notation $F_{ij}$ to denote the 
entries of $F$ in the $i$-th row and $j$-th column, i.e., $F_{ij} = f_{ij}$ in the above 
case. We will also use the shorthand notation $F = (f_{ij})$ to denote a matrix $F$ with entries $f_{ij}$ in the $i$-th row and $j$-th column.

With the above notation we can then define the matrix-vector product $F \xb$ as
\begin{equation*}
    F \xb = f(\xb) = \begin{bmatrix}
        \langle \fb_{1:}, \xb \rangle \\
        \langle \fb_{2:}, \xb \rangle \\
        \vdots \\
        \langle \fb_{m:}, \xb \rangle  
    \end{bmatrix} 
    =  
    \sum_{j=1}^n x_j \fb_{:j}.
\end{equation*}
In words, $F \xb$ denotes the evaluation of the underlying linear function $f$ at 
$\xb$ which can be represented in two ways: (1) as a list of inner products of the rows of $F$ with $\xb$ or (2) as a linear combination of the columns of $F$ with coefficients given by the entries of $\xb$.

For two matrices $F \in \R^{m \times n}$ and $G \in \R^{n \times p}$, we can define their product $FG \in \R^{m \times p}$ as the composition of the underlying linear functions, i.e., $FG \xb = f(g(\xb))$ for all $\xb \in \R^p$. In terms of the entries of $F$ and $G$, we have
\begin{equation*}
    (FG)_{ij} = \sum_{k=1}^n f_{ik} g_{kj} = \langle \fb_{i:}, \gb_{:j} \rangle.
\end{equation*}
Given $F \in \R^{m \times n}$, we can also define the transpose of $F$, denoted by $F^\top \in \R^{n \times m}$, as the matrix with entries $(F^\top)_{ij} = f_{ji}$ for all $i=1,2,\ldots,n$ and $j=1,2,\ldots,m$. In other words, the rows of $F$ become the columns of $F^\top$ and vice versa. Equivalently, we can write $F^\top = (f_{ji})$ to denote the transpose of $F$. The transpose $F^\top$ is the unique matrix in $\R^{n \times m}$ such that $\langle F \xb, \yb \rangle = \langle \xb, F^\top \yb \rangle$ for all $\xb \in \R^n$ and $\yb \in \R^m$.

\begin{definition}
    We say that a matrix $F$ is square if it has the same number of rows and columns, i.e., $F \in \R^{n \times n}$ for some $n \in \N$. We say that a square matrix $F$ is 
    symmetric if $F = F^\top$, i.e., $f_{ij} = f_{ji}$ for all $i,j=1,2,\ldots,n$.
\end{definition}

\begin{definition}
    A square matrix $F \in \R^{n \times n}$ is invertible if there exists a square matrix $G \in \R^{n \times n}$ such that $FG = GF = I_n$ where $I_n$ is the $n \times n$ identity matrix. The matrix $G$ is called the inverse of $F$ and is denoted by $F^{-1}$.
\end{definition}

\begin{lemma}
    Consider $F \in \R^{n \times n}$. Then the following statements are equivalent:

    \begin{itemize}
        \item $F$ is invertible.
        \item The columns of $F$ are linearly independent.
        \item The rows of $F$ are linearly independent.
        \item The equation $F \xb = \yb$ has a unique solution for every $\yb \in \R^n$.
        \item The equation $F \xb = 0$ has only the trivial solution $\xb = 0$.
    \end{itemize}
\end{lemma}

\begin{definition}
    We say that a square matrix $F \in \R^{n \times n}$ is positive definite if $\langle F \xb, \xb \rangle > 0$ for all $\xb \in \R^n$ with $\xb \neq 0$. We say that $F$ is positive semi-definite if $\langle F \xb, \xb \rangle \ge 0$ for all $\xb \in \R^n$.
\end{definition}

Positive definite matrices are some of the nicest objects in linear algebra. For example, 
they are always invertible and their inverses are also positive definite. Additionally, 
a positive definite matrix $F$ induces an inner product and norm on $\R^n$ defined as $\langle \xb, \yb \rangle_F := \langle F \xb, \yb \rangle$ and $\|\xb\|_F := \sqrt{\langle F \xb, \xb \rangle}$ for all $\xb, \yb \in \R^n$.

\section{Extension to complex variables}

All of the ideas and definitions above can be extended to complex variables. We denote the $d$-dimensional complex space by $\C^d$. For $\zb = (z_1, z_2, \ldots, z_d) \in \C^d$ we define its complex conjugate as $\overline{\zb} = (\overline{z_1}, \overline{z_2}, \ldots, \overline{z_d})$ where $\overline{z_j}$ is the complex conjugate of $z_j$ for $j=1,2,\ldots,d$. The Euclidean inner product on $\C^d$ is then defined as 
\begin{equation*}
    \langle \zb, \wb \rangle := \sum_{j=1}^d z_j \overline{w_j} = \langle \overline{\wb}, \overline{\zb} \rangle
\end{equation*}
The notions of norm, orthogonality, linear independence, span, basis, and dimension can be defined in the same way as in the real case. The Gram-Schmidt process can also be applied to complex vectors to obtain an orthonormal set of vectors just as before simply by 
changing our definition of the inner product to the complex inner product defined above.

We can also define complex matrices as before and denote their space as $\C^{n \times m}$. The transpose of a complex matrix $F \in \C^{n \times m}$ is defined as before and is denoted by $F^\top \in \C^{m \times n}$. Additionally, we can define the conjugate transpose of $F$, denoted by $F^* \in \C^{m \times n}$, as the matrix with entries $(F^*)_{ij} = \overline{f_{ji}}$ for all $i=1,2,\ldots,m$ and $j=1,2,\ldots,n$. In other words, the conjugate transpose of a complex matrix is obtained by taking the transpose and then taking the complex conjugate of each entry. The conjugate transpose is also sometimes called the Hermitian transpose and may be denoted as $F^H$ but we will not use this notation. Note that the conjugate transpose of a matrix 
is, in some sense, the more natural extension of the transpose to complex matrices since it satisfies the property $\langle F \zb, \wb \rangle = \langle \zb, F^* \wb \rangle$ for all $\zb \in \C^n$ and $\wb \in \C^m$.

Regardless, a symmetric complex matrix $F \in \R^{n \times n}$ is defined as a matrix that satisfies $F = F^\top$ and a Hermitian complex matrix $F \in \C^{n \times n}$ is defined as a matrix that satisfies $F = F^*$. A Hermitian matrix is the natural extension of a symmetric matrix to complex matrices. A Hermitian matrix $F$ is positive definite if $\langle F \zb, \zb \rangle > 0$ for all $\zb \in \C^n$ with $\zb \neq 0$ and it is positive semi-definite if $\langle F \zb, \zb \rangle \ge 0$ for all $\zb \in \C^n$.


\section{Some special types of matrices}

Certain matrices with special structure appear so frequently in applications and numerical computations that they deserve their own names. We collect a few of the most important ones here.

\begin{definition}[Identity matrix]
    The identity matrix $I_n \in \R^{n \times n}$ is defined entrywise by
    \begin{equation*}
        (I_n)_{ij} = \begin{cases} 1 & i = j \\ 0 & i \neq j \end{cases}, \qquad i,j = 1,2,\ldots,n.
    \end{equation*}
\end{definition}
For example,
\begin{equation*}
    I_3 = \begin{bmatrix}
        1 & 0 & 0 \\
        0 & 1 & 0 \\
        0 & 0 & 1
    \end{bmatrix}.
\end{equation*}
As the name suggests, $I_n$ is the multiplicative identity for matrix multiplication: $I_n F = F I_n = F$ for every $F \in \R^{n \times n}$, and $I_n \xb = \xb$ for every $\xb \in \R^n$. We already used $I_n$ in this way when defining the inverse of a matrix above.

\begin{definition}[Diagonal matrix]
    A square matrix $D \in \R^{n \times n}$ is called diagonal if $(D)_{ij} = 0$ whenever $i \neq j$. We write $D = \mathrm{diag}(d_1, d_2, \ldots, d_n)$ where $d_i = (D)_{ii}$ for $i=1,2,\ldots,n$.
\end{definition}
For example,
\begin{equation*}
    D = \begin{bmatrix}
        2 & 0 & 0 \\
        0 & -1 & 0 \\
        0 & 0 & 5
    \end{bmatrix} = \mathrm{diag}(2,-1,5).
\end{equation*}
The identity matrix is the special case $I_n = \mathrm{diag}(1,1,\ldots,1)$. Diagonal matrices act on vectors by simple entrywise scaling, $(D\xb)_i = d_i x_i$, and the product of two diagonal matrices is again diagonal, obtained by multiplying corresponding diagonal entries. A diagonal matrix $D$ is invertible if and only if $d_i \neq 0$ for every $i$, in which case $D^{-1} = \mathrm{diag}(1/d_1, 1/d_2, \ldots, 1/d_n)$.

\begin{definition}[Triangular matrices]
    A square matrix $U \in \R^{n \times n}$ is called upper triangular if $(U)_{ij} = 0$ whenever $i > j$, and a square matrix $L \in \R^{n \times n}$ is called lower triangular if $(L)_{ij} = 0$ whenever $i < j$.
\end{definition}
For example,
\begin{equation*}
    U = \begin{bmatrix}
        1 & 2 & 3 \\
        0 & 4 & 5 \\
        0 & 0 & 6
    \end{bmatrix}, \qquad
    L = \begin{bmatrix}
        1 & 0 & 0 \\
        2 & 4 & 0 \\
        3 & 5 & 6
    \end{bmatrix}
\end{equation*}
are, respectively, upper and lower triangular. Note that $L = U^\top$ in this example; in general the transpose of an upper triangular matrix is lower triangular and vice versa. Diagonal matrices are simultaneously upper and lower triangular. The sum or product of two upper (lower) triangular matrices is again upper (lower) triangular, and a triangular matrix is invertible if and only if all of its diagonal entries are nonzero, in which case its inverse is triangular of the same type. Triangular matrices are important because linear systems $U\xb = \yb$ or $L\xb = \yb$ can be solved efficiently by back- or forward-substitution; this is the basis of the $LU$ decomposition discussed below.

\begin{definition}[Orthogonal matrix]
    A square matrix $Q \in \R^{n \times n}$ is called orthogonal if
    \begin{equation*}
        Q^\top Q = Q Q^\top = I_n,
    \end{equation*}
    or equivalently $Q^{-1} = Q^\top$.
\end{definition}
Equivalently, $Q$ is orthogonal precisely when its columns (or, equivalently, its rows) form an orthonormal set of vectors in $\R^n$ with respect to the Euclidean inner product \eqref{def:euclidean-inner-product}. A useful consequence is that orthogonal matrices preserve inner products and norms:
\begin{equation*}
    \langle Q\xb, Q\yb \rangle = \langle \xb, \yb \rangle, \qquad \|Q\xb\|_2 = \|\xb\|_2, \qquad \text{for all } \xb, \yb \in \R^n.
\end{equation*}
For example, the permutation matrix
\begin{equation*}
    Q_1 = \begin{bmatrix}
        0 & 1 & 0 \\
        0 & 0 & 1 \\
        1 & 0 & 0
    \end{bmatrix}
\end{equation*}
is orthogonal since its columns are just the standard basis vectors of $\R^3$ in a different order. Another example is the rotation matrix
\begin{equation*}
    Q_2 = \begin{bmatrix}
        \cos\theta & -\sin\theta & 0 \\
        \sin\theta & \cos\theta & 0 \\
        0 & 0 & 1
    \end{bmatrix},
\end{equation*}
which rotates $\R^3$ by an angle $\theta$ about the third coordinate axis; one can check directly that $Q_2^\top Q_2 = I_3$ for any choice of $\theta$.

\begin{definition}[Unitary matrix]
    A square matrix $U \in \C^{n \times n}$ is called unitary if
    \begin{equation*}
        U^* U = U U^* = I_n,
    \end{equation*}
    or equivalently $U^{-1} = U^*$.
\end{definition}
As with the orthogonal case, $U$ is unitary precisely when its columns (or rows) form an orthonormal set of vectors in $\C^n$ with respect to the complex inner product introduced above, and unitary matrices preserve the complex inner product and norm in the same way that orthogonal matrices do on $\R^n$. In fact, orthogonal matrices are exactly the real unitary matrices: when all entries of $U$ are real, the conjugate transpose $U^*$ coincides with the transpose $U^\top$. A simple example of a unitary matrix is a diagonal matrix whose entries have unit modulus, such as
\begin{equation*}
    U_1 = \begin{bmatrix}
        1 & 0 & 0 \\
        0 & i & 0 \\
        0 & 0 & -1
    \end{bmatrix},
\end{equation*}
since $U_1^* U_1 = \mathrm{diag}(|1|^2, |i|^2, |-1|^2) = I_3$. A less trivial example, which will reappear later in the course in connection with the discrete Fourier transform, is
\begin{equation*}
    U_2 = \frac{1}{\sqrt{3}}\begin{bmatrix}
        1 & 1 & 1 \\
        1 & \omega & \omega^2 \\
        1 & \omega^2 & \omega
    \end{bmatrix}, \qquad \omega = e^{2\pi i/3},
\end{equation*}
whose columns can be checked to be orthonormal with respect to the complex inner product.

\section{Useful matrix decompositions}

A matrix decomposition rewrites a matrix as a product of matrices with simpler structure, such as the special types introduced above. Decompositions are central to numerical linear algebra: they are the standard tools for solving linear systems, least-squares problems, and eigenvalue problems efficiently and stably. We simply state the definitions here; the algorithms used to compute these decompositions can be 
found in classic references such as \cite{demmel1997applied,olver2006applied}.

\begin{definition}[LU decomposition]
    Let $F \in \R^{n \times n}$ be a matrix that can be reduced to upper triangular form using only row operations that add a multiple of one row to another (i.e., without row interchanges). Then there exist a unit lower triangular matrix $L \in \R^{n \times n}$ (i.e., lower triangular with $(L)_{ii}=1$ for all $i$) and an upper triangular matrix $U \in \R^{n \times n}$ such that
    \begin{equation*}
        F = LU.
    \end{equation*}
\end{definition}
The entries of $L$ record the multipliers used during Gaussian elimination, and $U$ is the resulting upper triangular matrix. Once $F$ has been factored in this way, the linear system $F\xb = \yb$ reduces to two easy triangular solves: $L\zb = \yb$ followed by $U\xb = \zb$. For example,
\begin{equation*}
    F = \begin{bmatrix}
        2 & 1 & 1 \\
        4 & 3 & 3 \\
        8 & 7 & 9
    \end{bmatrix}
    =
    \underbrace{\begin{bmatrix}
        1 & 0 & 0 \\
        2 & 1 & 0 \\
        4 & 3 & 1
    \end{bmatrix}}_{L}
    \underbrace{\begin{bmatrix}
        2 & 1 & 1 \\
        0 & 1 & 1 \\
        0 & 0 & 2
    \end{bmatrix}}_{U}.
\end{equation*}

\begin{definition}[QR decomposition]
    Let $F \in \R^{n \times n}$ be invertible. Then there exist an orthogonal matrix $Q \in \R^{n \times n}$ and an upper triangular matrix $R \in \R^{n \times n}$ with positive diagonal entries such that
    \begin{equation*}
        F = QR.
    \end{equation*}
\end{definition}
This decomposition is precisely what results from applying the Gram-Schmidt process to the columns of $F$: the columns of $Q$ are the resulting orthonormal vectors, and the entries of $R$ record the coefficients used in forming each column of $F$ as a linear combination of the columns of $Q$. For example,
\begin{equation*}
    F = \begin{bmatrix}
        12 & -51 & 4 \\
        6 & 167 & -68 \\
        -4 & 24 & -41
    \end{bmatrix}
    =
    \underbrace{\frac{1}{175}\begin{bmatrix}
        150 & -69 & -58 \\
        75 & 158 & 6 \\
        -50 & 30 & -165
    \end{bmatrix}}_{Q}
    \underbrace{\begin{bmatrix}
        14 & 21 & -14 \\
        0 & 175 & -70 \\
        0 & 0 & 35
    \end{bmatrix}}_{R}.
\end{equation*}

\begin{definition}[Cholesky decomposition]
    Let $F \in \R^{n \times n}$ be symmetric positive definite. Then there exists a unique lower triangular matrix $L \in \R^{n \times n}$ with positive diagonal entries such that
    \begin{equation*}
        F = LL^\top.
    \end{equation*}
\end{definition}
The Cholesky decomposition can be viewed as a kind of matrix square root for positive definite matrices, and also as a symmetric special case of the $LU$ decomposition in which $U = L^\top$. For example,
\begin{equation*}
    F = \begin{bmatrix}
        4 & 12 & -16 \\
        12 & 37 & -43 \\
        -16 & -43 & 98
    \end{bmatrix}
    =
    \underbrace{\begin{bmatrix}
        2 & 0 & 0 \\
        6 & 1 & 0 \\
        -8 & 5 & 3
    \end{bmatrix}}_{L}
    \underbrace{\begin{bmatrix}
        2 & 6 & -8 \\
        0 & 1 & 5 \\
        0 & 0 & 3
    \end{bmatrix}}_{L^\top}.
\end{equation*}

\begin{definition}[Singular value decomposition]
    Let $F \in \R^{m \times n}$ be any matrix (not necessarily square). Then there exist orthogonal matrices $U \in \R^{m \times m}$ and $V \in \R^{n \times n}$, and a matrix $\Sigma \in \R^{m \times n}$ that is zero off of its main diagonal with nonnegative diagonal entries $\sigma_1 \ge \sigma_2 \ge \cdots \ge 0$, such that
    \begin{equation*}
        F = U \Sigma V^\top.
    \end{equation*}
    The $\sigma_j$ are called the singular values of $F$, and the columns of $U$ and $V$ are called, respectively, the left and right singular vectors of $F$.
\end{definition}
Unlike the previous three decompositions, the SVD exists for \emph{every} matrix $F$, regardless of its shape or rank, which makes it one of the most versatile tools in numerical linear algebra. For example,
\begin{equation*}
    F = \begin{bmatrix}
        4 & 1 & 6 \\
        -2 & -2 & 6 \\
        -4 & 2 & 3
    \end{bmatrix}
    =
    \underbrace{\frac{1}{3}\begin{bmatrix}
        2 & 2 & 1 \\
        2 & -1 & -2 \\
        1 & -2 & 2
    \end{bmatrix}}_{U}
    \underbrace{\begin{bmatrix}
        9 & 0 & 0 \\
        0 & 6 & 0 \\
        0 & 0 & 3
    \end{bmatrix}}_{\Sigma}
    \underbrace{\begin{bmatrix}
        0 & 0 & 1 \\
        1 & 0 & 0 \\
        0 & 1 & 0
    \end{bmatrix}}_{V^\top}.
\end{equation*}


\section{Eigenvalue decompositions}

Throughout this course we will heavily rely on the eigenvalue decomposition of 
a matrix to make sense of many ideas from stability and convergence 
properties of ODEs and numerical algorithms to the design of 
efficient algorithms such as spectral methods. 

\begin{definition}
    Consider $F \in \C^{n \times n}$. A nonzero vector $\qb \in \C^n$ is called an eigenvector of $F$ if there exists a scalar $\lambda \in \C$ such that
    \begin{equation*}
        F \qb = \lambda \qb.
    \end{equation*}
    We call $(\lambda, \qb)$ an eigenpair of $F$.
\end{definition}
We say that $F$ is diagonalizable if it has $n$ 
linearly independent eigenvectors. In that case, we obtain an eigenvalue 
decomposition (eigen-decomposition for short) of $F$ as
\begin{equation*}
    F = Q \Lambda Q^{-1}
\end{equation*}
where $Q \in \C^{n \times n}$ is the matrix whose columns are
the eigenvectors of $F$ and $\Lambda \in \C^{n \times n}$ is the diagonal matrix whose diagonal entries are the corresponding eigenvalues of $F$.

It is important to note that not every matrix is diagonalizable 
and the eigendecomposition does not exist like the SVD.

\begin{example}
    The following matrix is not diagonalizable:
    \begin{equation*}
        F 
= \begin{bmatrix}
            1 & 1 \\
            0 & 1
        \end{bmatrix}
    \end{equation*}
To see this try to compute the eigenvalues of $F$ by solving the characteristic polynomial $\det(F - \lambda I) = 0$. We have
\begin{equation*}
    \det(F - \lambda I) = \det\begin{bmatrix}
        1 - \lambda & 1 \\
        0 & 1 - \lambda
    \end{bmatrix} = (1 - \lambda)^2.
\end{equation*}
Hence, the only eigenvalue of $F$ is $\lambda = 1$. To find the corresponding eigenvectors we solve the equation $(F - I)\qb = 0$ which gives
\begin{equation*}
    \begin{bmatrix}
        0 & 1 \\
        0 & 0
    \end{bmatrix} \begin{bmatrix}
        q_1 \\
        q_2
    \end{bmatrix} = \begin{bmatrix}
        0 \\
        0
    \end{bmatrix}.
\end{equation*}
This gives the equation $q_2 = 0$ and hence the eigenvectors of $F$ are of the form $\qb = (q_1, 0)^\top$. In particular, there is only one linearly independent eigenvector of $F$ and hence $F$ is not diagonalizable.
\end{example}

\begin{lemma}
    Suppose $F \in \C^{n \times n}$ is Hermitian. 
    Then $F$ is diagonalizable and all of its eigenvalues are real, i.e., 
    there exists a unitary matrix $Q \in \C^{n \times n}$ and a diagonal matrix $\Lambda \in \R^{n \times n}$ such that
    \begin{equation*}
        F = Q \Lambda Q^*.
    \end{equation*}
    If $F$ is also positive semi-definite, then all of its eigenvalues are nonnegative. If $F$ is positive definite, then all of its eigenvalues are positive.
\end{lemma}

\begin{example}
    Consider the matrix 
    \begin{equation*}
        F = \begin{bmatrix}
            1 & - i & 0 \\ 
            i & 1 & -i \\ 
            0 & i & 1
        \end{bmatrix}
    \end{equation*}
    It is easy to check that $F$ is Hermitian since $F^* = F$. Then it has an eigen-decomposition $F = Q \Lambda Q^*$ where
    \begin{equation*}
        Q =  \begin{bmatrix}
            \frac12 & \frac{1}{\sqrt{2}} & \frac12 \\
            -\frac{i}{\sqrt{2}} & 0 & \frac{i}{\sqrt{2}} \\
            - \frac12 & \frac{1}{\sqrt{2}} & -\frac12
        \end{bmatrix}, \qquad
        \Lambda = \begin{bmatrix}
            1 - \sqrt{2} & 0 & 0 \\
            0 & 1 & 0 \\
            0 & 0 & 1 + \sqrt{2}
        \end{bmatrix}.
    \end{equation*}
\end{example}


\section{Matrix powers and exponential}
For a square matrix $F \in \C^{n \times n}$ we can define its powers $F^k$ for $k \in \N$ as the product of $k$ copies of $F$, i.e., $F^k = F F \cdots F$ ($k$ times). We also define $F^0 = I_n$. If $F$ is diagonalizable with eigen-decomposition $F = Q \Lambda Q^{-1}$, then we can compute its powers as
\begin{equation*}
    F^k = Q \Lambda^k Q^{-1}
\end{equation*}
where $\Lambda^k$ is the diagonal matrix with entries $(\Lambda^k)_{ii} = \lambda_i^k$ for $i=1,2,\ldots,n$. This is a very useful property since it allows us to compute powers of $F$ by simply raising its eigenvalues to the $k$-th power.

With matrix powers at hand, we can define polynomial functions of 
matrices $p(F)$ for a polynomial $p(x) = \sum_{k=0}^m a_k x^k$ as
\begin{equation*}
    p(F) = \sum_{k=0}^m a_k F^k.
\end{equation*}
 If $F$ is diagonalizable this takes the form
\begin{equation*}
p(F) = Q p(\Lambda) Q^{-1}
\end{equation*}
where $p(\Lambda)$ is the diagonal matrix with entries $(p(\Lambda))_{ii} = p(\lambda_i)$ for $i=1,2,\ldots,n$. 

Finally, using the Taylor series expansion of the exponential function, we can define the matrix exponential as
\begin{equation*}
    e^F = \sum_{k=0}^\infty \frac{F^k}{k!},
\end{equation*}
which for a diagonalizable matrix $F$ takes the simpler form 
\begin{equation*}
    e^F = Q e^{\Lambda} Q^{-1}.
\end{equation*}

\section{Kronecker products}

The Kronecker product of two matrices $A \in \C^{m \times n}$ and $B \in \C^{p \times q}$ is defined as the block matrix
\begin{equation*}
    A \otimes B = \begin{bmatrix}
        a_{11} B & a_{12} B & \cdots & a_{1n} B \\
        a_{21} B & a_{22} B & \cdots & a_{2n} B \\
        \vdots & \vdots & \ddots & \vdots \\
        a_{m1} B & a_{m2} B & \cdots & a_{mn} B
    \end{bmatrix} \in \C^{mp \times nq}.
\end{equation*}
Some important properties of the Kronecker product include:
\begin{itemize}
    \item $(A \otimes B)(C \otimes D) = (AC) \otimes (BD)$, provided the matrix multiplications $AC$ and $BD$ are defined.
    \item $(A \otimes B)^* = A^* \otimes B^*$.
    \item If $A$ and $B$ are invertible, then $(A \otimes B)^{-1} = A^{-1} \otimes B^{-1}$.
    \item If $A$ and $B$ are diagonalizable with $A = Q_A \Lambda_A Q_A^{-1}$ and $B = Q_B \Lambda_B Q_B^{-1}$, then
    \begin{equation*}
        A \otimes B = (Q_A \otimes Q_B)(\Lambda_A \otimes \Lambda_B)(Q_A \otimes Q_B)^{-1}.
    \end{equation*}
\end{itemize}
Perhaps most consequential for us is the property that the Kronecker product interacts nicely with the vectorization of matrices. Specifically, for matrices $A \in \C^{m \times n}$, $B \in \C^{p \times q}$, and $X \in \C^{n \times p}$, we have
\begin{equation}\label{eq:vec-kronecker}
    \text{vec}(AXB) = (B^T \otimes A) \text{vec}(X),
\end{equation}
where $\text{vec}(X)$ denotes the vectorization of $X$ obtained by stacking its columns into a single column vector.


\bibliographystyle{plainnat}
\bibliography{references}





